\documentclass[11pt]{article}
\usepackage[utf8]{inputenc}
\usepackage[T1]{fontenc}
\usepackage[margin=2.5cm]{geometry}
\usepackage{hyperref}
\title{Errata for the cloud Claude pass (ranges t--tzz and u--uzz)}
\author{Compiled by Claude Sonnet 5 (claude-sonnet-5) for Andr\'as Kornai}
\date{14 October 2026}
\begin{document}
\maketitle
\noindent Dossiers were saved as the subagents returned them; the items below were noticed
during the run or found by a post-hoc scan of \texttt{K12a/[tu]??}. Nothing here has been
corrected in place unless stated. The generated LaTeX was never compiled (no TeX in the
cloud container), so the compile warnings are suspicions, not test results.

\section{Corrected before saving}
\begin{itemize}
\item \textbf{usp}: the subagent had written \texttt{glorrolog.invalid} as a URL; removed, citation (Can Pixabaj 2007, \emph{Gram\'atica Uspanteka}) kept as plain text.
\item \textbf{tac}: accents altered in transit; restored.
\item \textbf{tew, thk}: the JSON block was invalid; \texttt{resources\_sonnet.json} was built from the URLs in the LaTeX block instead.
\end{itemize}

\section{Likely wrong or mismatched source}
\begin{itemize}
\item \textbf{umb}: the dbs.org UBUAND page may describe Umbu-Ungu, a different language.
\item \textbf{uuu}: the Glottolog (\texttt{mang1378}) and ELP (entry 1236) links may be another language; the dossier says so.
\item \textbf{uwa}: the prompt calls Kuku-Uwanh extinct, the sources describe a dialect of Kugu-Nganhcara with a few elderly speakers; unresolved. Shares an ELP page (node/92026) with \textbf{ugb}.
\item \textbf{usu, utu}: both cite ELP \texttt{lang/3409}; at least one is wrong. ``Uya'' is also a name of Sausi (ssj).
\item \textbf{tvd, tyy}: same ELP page (node/96036). \textbf{udj, urn}: same ELP page (lang/2798). \textbf{tbn, tnd}: same ELP page (lang/3527) and Tunebo PDF.
\item \textbf{uma}: a CTUIR language program contradicts the prompt's ``not known to be taught''.
\item \textbf{unx} is a retired code (same language as Mundari, \texttt{unr}); \textbf{uok} is doubtful.
\item \textbf{uzn}: one Dunyo.info URL is percent-garbled and probably truncated. \textbf{uvl}: a Zenodo DOI appears as plain text and was not seen in search results. \textbf{ute}: only Southern Ute was searched, not Southern Paiute.
\item The family field is wrong or bookkeeping-only in several dossiers.
\end{itemize}

\section{LaTeX that may not compile}
\textbf{tyt} (\verb|\uhorn|, \verb|\ohorn|), \textbf{tzl} (\verb|\th{}|), \textbf{urd} (a raw space inside a \verb|\url|), \textbf{uuu} (\verb|\h{a}|).

\section{URLs shared between dossiers (scan, 14 Oct 2026)}
Most are legitimate shared sources for related languages (family surveys, Webonary
Tamang, Tanna bibliographies, Wikipedia). Pairs worth a check beyond those above:
tac/twr, tlf/urm, tkw/tnx, tze/tzz, thz/ttq, urp/urz (Lingua Viva article), uzb/uzn
(expected: macrolanguage and member). The Tanna pages are shared by tnk, tnl, tnn, tnp.

\end{document}
